\chapter{Considerações Finais e Plano de Trabalho}
\label{chap:consideracoes_plano}

Após a exposição dos fundamentos teóricos, a análise crítica do estado da arte e a proposição detalhada da arquitetura e do mapeamento gesto–som, este capítulo consolida as contribuições alcançadas nesta etapa do trabalho e apresenta o planejamento para a fase de implementação, avaliação e documentação do mesmo. Inicialmente, resumem-se as conclusões principais e, em seguida, descreve-se de forma objetiva o que se espera entregar no TCC 2, detalhando o fluxo de atividades e o cronograma de execução proposto.

\section{Considerações Finais do TCC 1}
\label{sec:conclusao_tcc1}

Nesta etapa inicial, foi identificado o problema central de barreiras de acesso à expressão musical, especialmente em cenários de interfaces gestuais de baixo custo e alta expressividade. A análise do contexto tecnológico revelou que, embora existam soluções pontuais, persistia uma lacuna: a falta de um sistema que conjunte alta acessibilidade, baixa latência e alta expressividade. Essa constatação orientou o levantamento bibliográfico e a análise crítica apresentada no Capítulo~\ref{cap:analise_estado_arte_interfaces_musicais_gestuais}, a qual reafirmou que não há, no estado da arte, proposta que equilibre esses três pilares de forma integrada.

Em resposta a essa lacuna, o projeto Aerochord propõe uma solução: um hiperinstrumento digital orientado a gestos que utiliza apenas uma câmera RGB e software de baixo nível (C++/JUCE) para capturar, mapear e gerar eventos MIDI 2.0 de alta resolução. A principal contribuição do TCC 1 consiste na entrega de uma proposta de engenharia completa e bem fundamentada, incluindo uma revisão extensiva do referencial teórico e do estado da arte em interfaces gestuais musicais, identificação clara das limitações das abordagens existentes, definição de requisitos funcionais e não funcionais, detalhamento de uma arquitetura modular para implementação futura e proposição de mapeamentos gesto–som simples e coerentes com as necessidades de usabilidade. Este conjunto de artefatos forma a base sólida para a fase de desenvolvimento e avaliação empiricamente orientada no TCC 2.

\section{Resultados Esperados para o TCC 2}
\label{sec:resultados_esperados_tcc2}

Espera-se que a próxima etapa do projeto (TCC 2) entregue os seguintes artefatos e resultados, os quais deverão ser desenvolvidos e validados conforme descrito:

\begin{itemize}
    \item \textbf{Protótipo de Software (concluído):} Um protótipo funcional do hiperinstrumento Aerochord, implementado em C++/JUCE, que realiza o \textit{pipeline} completo de captura de vídeo, detecção de pose via MediaPipe, análise de gestos, mapeamento gesto–som e geração de mensagens MIDI 2.0 (com \textit{fallback} para MIDI 1.0), conforme a arquitetura descrita no Capítulo~\ref{chap:desenvolvimento_proposta}. Esta etapa avançou além do previsto para o início do TCC 2: o núcleo do sistema, a lógica de interação e a integração entre os módulos já foram implementados e validados por testes automatizados.
    \item \textbf{Validação Experimental:} A condução de protocolo formal de avaliação com usuários-alvo, incluindo medições objetivas (latência \textit{end-to-end} medida via \textit{timestamps} em UMP, taxa de reconhecimento de gestos, estabilidade do rastreamento) e subjetivas (questionários SUS para usabilidade inicial, NASA-TLX para carga cognitiva, entrevistas semiestruturadas para \textit{feedback} qualitativo). Os resultados devem embasar refinamentos no mapeamento e ajustes de parâmetros sensíveis.
    \item \textbf{Monografia Final:} A redação dos capítulos restantes da dissertação, englobando Metodologia de Desenvolvimento, Descrição de Implementação, Resultados da Avaliação, Discussão dos Resultados, Conclusões e Trabalhos Futuros, formando um documento coeso que apresente toda a jornada do projeto Aerochord desde a ideia até a validação.
\end{itemize}

\section{Plano de Atividades para o TCC 2}
\label{sec:plano_atividades_tcc2}

Para estruturar o trabalho a ser realizado, apresenta-se a seguir um plano de atividades dividido em fases sequenciais, com estimativa de duração em semanas. Cada fase foca em um conjunto de tarefas articuladas para garantir progresso incremental e avaliação contínua. As três primeiras fases, referentes ao desenvolvimento do núcleo do sistema, já foram concluídas durante a elaboração deste documento, adiantando-se em relação ao cronograma originalmente previsto.

\begin{enumerate}
    \item \textbf{Fase 1 – Desenvolvimento do Núcleo do Sistema (Semanas 1–5) — Concluída:}
    Foram implementados os módulos básicos: captura de vídeo via V4L2 com \textit{buffering} circular, integração do MediaPipe Hand Landmarker para extração de \textit{landmarks} de mãos (com filtro EMA de suavização), e módulo de geração de mensagens MIDI 2.0 em C++/JUCE, incluindo \textit{handshake} MIDI-CI via ALSA Sequencer e \textit{fallback} automático para MIDI 1.0. As interfaces internas (filas \textit{lock-free} entre módulos) foram validadas e o envio de pacotes UMP foi testado de ponta a ponta.
    \item \textbf{Fase 2 – Implementação da Lógica de Interação (Semanas 4–7) — Concluída:}
    Foi desenvolvido o módulo de análise de gestos, baseado em máquina de estados finitos, para reconhecimento de gestos de pinça, seleção de oitava e controles contínuos (\textit{pitch bend}, \textit{timbre}, \textit{vibrato}, \textit{sustain}). Em paralelo, implementou-se o módulo de mapeamento gesto–som, conforme a Tabela~\ref{tab:mapeamento_implementado} e a Seção~\ref{sec:estrategias_mapeamento} do Capítulo~\ref{chap:desenvolvimento_proposta}. Mecanismos de \textit{debounce}, histerese e filtros de suavização nos \textit{landmarks} foram aplicados para mitigar falsos positivos.
    \item \textbf{Fase 3 – Integração e Testes Iniciais (Semanas 8–10) — Concluída:}
    Todos os módulos implementados foram integrados em um \textit{pipeline} único, com comunicação por filas \textit{lock-free} de ponta a ponta. Testes automatizados (unitários e de integração da cadeia completa gesto–MIDI) confirmam que captura, detecção, análise de gestos, mapeamento e geração MIDI operam em conjunto sem travamentos. Os ajustes de parâmetros (\textit{buffers}, \textit{thresholds} de confiança, zonas de mapeamento) seguem sendo refinados com base em medições de \textit{profiling}.
    \item \textbf{Fase 4 – Execução da Avaliação com Usuários (Semanas 11–12):}
    Planeja-se recrutar participantes representativos do público-alvo (músicos sem treinamento formal e pessoas com mobilidade reduzida) e conduzir sessões de uso do protótipo em ambiente controlado. Será aplicado protocolo de avaliação conforme metodologia discutida na Seção~\ref{sec:avaliacao_usuarios_contextos}, coletando dados objetivos (latência, precisão de gestos) e subjetivos (SUS, NASA-TLX, entrevistas). Os logs de eventos gestuais e de mensagens MIDI serão registrados para análise posterior.
    \item \textbf{Fase 5 – Análise dos Dados e Escrita dos Resultados (Semanas 13–14):}
    Com os dados coletados, realizará-se análise estatística e qualitativa, comparando métricas de desempenho com metas estabelecidas (por exemplo, latência < 30 ms). Os resultados serão discutidos em relação ao referencial teórico e às expectativas de expressividade e usabilidade. Em seguida, redigirá-se o capítulo de Resultados e Discussão, destacando aprendizados e ajustes de mapeamento derivados da avaliação.
    \item \textbf{Fase 6 – Redação Final e Preparação para Defesa (Semanas 15–16):}
    Inclui a revisão completa da monografia, incorporação de \textit{feedback} do orientador, finalização de diagramas e tabelas, e elaboração de apresentação para a banca. Será realizada verificação de formatação e consistência de citações, além de planejamento da defesa para comunicar de forma clara a contribuição do Aerochord.
\end{enumerate}

\section{Cronograma de Execução Proposto}
\label{sec:cronograma}

Para demonstrar a viabilidade temporal do projeto, apresenta-se a seguir um cronograma simplificado distribuído ao longo de um semestre letivo de quatro meses (aproximadamente 16 semanas). Cada fase definida anteriormente é mapeada para os meses correspondentes, de modo a fornecer uma visão geral de alocação de tempo.

\begin{table}[htbp]
    \centering
    \caption{Cronograma Simplificado de Atividades do TCC 2}
    \label{tab:cronograma_tcc2}
    \scriptsize
    \setlength{\tabcolsep}{4pt}
    \renewcommand{\arraystretch}{1.0}
    \begin{tabular}{|l|c|c|c|c|}
        \hline
        \textbf{Atividade / Fase}&\textbf{Mês 1}&\textbf{Mês 2} & \textbf{Mês 3}             & \textbf{Mês 4}             \\ \hline
        Fase 1: Desenvolvimento do Núcleo do Sistema       & \textbullet & \textbullet &                            &                            \\ \cline{2-5} 
        Fase 2: Implementação da Lógica de Interação       &                            & \textbullet & \textbullet &                            \\ \cline{2-5} 
        Fase 3: Integração e Testes Iniciais               &                            &                            & \textbullet &                            \\ \cline{2-5} 
        Fase 4: Execução da Avaliação com Usuários         &                            &                            & \textbullet & \textbullet \\ \cline{2-5} 
        Fase 5: Análise dos Dados e Escrita dos Resultados &                            &                            &                            & \textbullet \\ \cline{2-5} 
        Fase 6: Redação Final e Preparação para Defesa     &                            &                            &                            & \textbullet \\ \hline
    \end{tabular}
\end{table}

O cronograma da tabela~\ref{tab:cronograma_tcc2} oferece uma visão geral de como as fases se sobrepõem ao longo dos meses: o desenvolvimento do núcleo inicia no Mês 1 e se estende até parte do Mês 2, a lógica de interação ocupa principalmente o Mês 2 e início do Mês 3, e assim por diante. Durante a execução, pode-se detalhar ainda mais por semana, mas este formato demonstra que as atividades planejadas são compatíveis com o período disponível e permitem margens para ajustes e imprevistos.  

Com isso, conclui-se o planejamento para o TCC 2, estabelecendo metas claras, entregáveis definidos e cronograma viável, sustentado pela fundamentação deste TCC 1 e pelas metodologias de avaliação e engenharia propostas nos capítulos anteriores.